\documentclass{article}
\usepackage[UTF8]{ctex}
\usepackage{amsmath,amssymb,amsfonts}
\usepackage{enumitem}
\usepackage{geometry}
\geometry{a4paper, margin=1in}
% === 错题本排版优化（不参与题目标识） ===
\setlength{\parindent}{0pt}
\setlength{\parskip}{0.35em plus 0.10em minus 0.08em}
\setlength{\emergencystretch}{3em}
\setlist{itemsep=0.25em, topsep=0.40em, parsep=0pt}
\allowdisplaybreaks[2]
% === 错题本排版优化结束 ===
\begin{document}
\section*{重庆统考2026年数学分析真题}
\section*{61．重庆统考2026年数学分析试卷}

\subsection{题目1}
1．求极限 \(\lim _{x \rightarrow+\infty} \frac{\ln \sqrt{x}+\sin ^{2} x}{x}\) ．

\subsection{题目2}
2．设 \(a>0\) ，求极限 \(\lim _{n \rightarrow \infty} \sqrt{a} \sqrt[4]{a} \sqrt[8]{a} \cdots \sqrt[2 n]{a}\) 。

\subsection{题目3}
3．求极限 \(\lim _{x \rightarrow+\infty}\left[x-x^{3} \ln \left(1+\frac{1}{x^{2}}\right)\right]\) ．

\subsection{题目4}
4．求极限

\[
\lim _{n \rightarrow \infty}\left(\frac{1}{\sqrt{n^{2}+1^{2}}}+\frac{1}{\sqrt{n^{2}+2^{2}}}+\cdots+\frac{1}{\sqrt{n^{2}+n^{2}}}\right) .
\]

\subsection{题目5}
5．已知函数

\[
g(t)= \begin{cases}t+2, & t \geq 0 \\ A t^{3}+B t^{2}+C t+D, & -r<t<0 \\ 0, & t \leq-r\end{cases}
\]

连续可微，其中 \(r>0\) ．

\begin{itemize}
  \item[（1）] 求 \(A, B, C, D\) 的值．
  \item[（2）] 求导函数 \(g^{\prime}(t)\) ．
\end{itemize}

\subsection{题目6}
6．已知函数 \(f(x)=\int_{1}^{x^{2}} \frac{\sin t}{t} \mathrm{~d} t\) ，求 \(\int_{0}^{1} x(f(x)+1) \mathrm{d} x\) ．

\subsection{题目7}
7．求幂级数 \(\sum_{n=1}^{+\infty} \frac{x^{n}}{n+1}\) 的收敛域与和函数。

\subsection{题目8}
8 ．求由方程 \(z \sin y+x e^{z}=1\) 所确定的隐函数 \(z=f(x, y)\) 在点 \((1,0)\) 的全微分．

\subsection{题目9}
9．计算曲线积分 \(\oint_{L} x e^{2 x y} \mathrm{~d} s\) ，其中 \(L\) 是以 \(O(0,0), A(1,0), B(1,1)\) 为顶点的三角形．

\subsection{题目10}
10．求曲面积分

\[
\iint_{S} x \mathrm{~d} y \mathrm{~d} z+y \mathrm{~d} z \mathrm{~d} x+z \mathrm{~d} x \mathrm{~d} y .
\]

其中 \(S\) 为曲面 \(z=\sqrt{x^{2}+y^{2}}\) 与 \(x^{2}+y^{2}+z^{2}=2\) 所围立体的表面，方向取外侧．

\subsection{题目11}
11．设 \(F(\theta)=\int_{0}^{\pi} \ln (1-\theta \cos x) \mathrm{d} x, \theta \in(-1,1)\) ．

\begin{itemize}
  \item[（1）] 求 \(F^{\prime}(\theta)\) 。
  \item[（2）] 求 \(F\left(\frac{\sqrt{3}}{2}\right)\) 的值．
\end{itemize}

\subsection{题目12}
12．点 \(M_{1}\left(x_{1}, y_{1}\right), M_{2}\left(x_{2}, y_{2}\right)\) 的距离定义为

\[
\begin{gathered}
d_{1}\left(M_{1}, M_{2}\right)=\sqrt{\left|x_{2}-x_{1}\right|^{2}+\left(y_{2}-y_{1}\right)^{2}} \\
d_{2}\left(M_{1}, M_{2}\right)=\left|x_{2}-x_{1}\right|+\left|y_{2}-y_{1}\right|
\end{gathered}
\]

证明：平面点列 \(\left\{\left(x_{n}, y_{n}\right)\right\}\) 按距离 \(d_{1}\) 收敛到 \(\left(x_{0}, y_{0}\right)\) 的充要条件是 \(\left\{\left(x_{n}, y_{n}\right)\right\}\) 按距离 \(d_{2}\) 收敛到\(\left(x_{0}, y_{0}\right)\) ．

\subsection{题目13}
13．设 \(a>1\) ，证明不等式： \(2 a^{\frac{1}{2}} \leq a^{x}+a^{1-x} \leq 1+a, x \in[0,1]\) ．

\subsection{题目14}
14．已知

\[
f(x, y)= \begin{cases}\left(x^{2}+y^{2}\right) \cos \frac{1}{x^{2}+y^{2}}, & x^{2}+y^{2} \neq 0 \\ 0, & x^{2}+y^{2}=0\end{cases}
\]

\begin{itemize}
  \item[（1）] 证明：函数 \(f(x, y)\) 在 \((0,0)\) 处可微．
  \item[（2）] 证明：函数 \(f(x, y)\) 的偏导数在 \((0,0)\) 处不连续．
\end{itemize}

\subsection{题目15}
15．已知 \(f(x)\) 在 \((a,+\infty)\) 上可导，且 \(\lim _{x \rightarrow+\infty} f^{\prime}(x)=A\) ．

\begin{itemize}
  \item[（1）] 若 \(A>0\) ，证明： \(\lim _{x \rightarrow+\infty} f(x)=+\infty\) ．
  \item[（2）] 若 \(A=0\) ，证明： \(\lim _{x \rightarrow+\infty} \frac{f(x)}{x}=0\) ．
\end{itemize}
\end{document}
